% AY2022 力学 第2問 転記（問題のみ）
\section*{第2問〔II〕}

図2のように，半径 $a$，慣性モーメント $I$，質量 $M$ の一様な円板の外周に太さの無視
できる軽い糸を巻き付け，糸の一端を天井に固定して静止させたのち，静かに手を離した
ところ，円板は，重心 G を通る垂直な軸のまわりを回転しながら，鉛直下向きに落下した。
このとき，次の問い (1)〜(3) に答えよ。ここで，重力加速度の大きさを $g$ とする。

\begin{enumerate}[label=(\arabic*)]
  \item 円板の角速度を $\omega$ として，$\left|\dfrac{d\omega}{dt}\right|$ を求めよ。
  \item 円板の加速度の大きさを求めよ。
  \item 糸の張力の大きさを求めよ。
\end{enumerate}

\begin{center}
% 図2：天井に固定した糸を外周に巻いた円板が落下。重心G、半径a
\begin{tikzpicture}[>=stealth, scale=1.0]
  \fill[pattern=north east lines] (-2.2,2.0) rectangle (2.2,2.16);
  \draw (-2.2,2.0) -- (2.2,2.0);
  \draw (0,2.0) -- (0,1.2);                         % 糸
  \draw (0,0.2) circle (1.0);                        % 円板
  \fill (0,0.2) circle (1.3pt); \node[right] at (0.08,0.2) {G};
  \draw[dashed] (0,0.2) -- (-1.0,0.2); \node[above] at (-0.55,0.2) {$a$};
  \node at (1.4,-0.9) {図 2};
\end{tikzpicture}
\end{center}
